% Options for packages loaded elsewhere
\PassOptionsToPackage{unicode}{hyperref}
\PassOptionsToPackage{hyphens}{url}
\documentclass[
]{article}
\usepackage{xcolor}
\usepackage[margin=1in]{geometry}
\usepackage{amsmath,amssymb}
\setcounter{secnumdepth}{-\maxdimen} % remove section numbering
\usepackage{iftex}
\ifPDFTeX
  \usepackage[T1]{fontenc}
  \usepackage[utf8]{inputenc}
  \usepackage{textcomp} % provide euro and other symbols
\else % if luatex or xetex
  \usepackage{unicode-math} % this also loads fontspec
  \defaultfontfeatures{Scale=MatchLowercase}
  \defaultfontfeatures[\rmfamily]{Ligatures=TeX,Scale=1}
\fi
\usepackage[]{charter}
\ifPDFTeX\else
  % xetex/luatex font selection
  \setmathfont[]{charter}
\fi
% Use upquote if available, for straight quotes in verbatim environments
\IfFileExists{upquote.sty}{\usepackage{upquote}}{}
\IfFileExists{microtype.sty}{% use microtype if available
  \usepackage[]{microtype}
  \UseMicrotypeSet[protrusion]{basicmath} % disable protrusion for tt fonts
}{}
\makeatletter
\@ifundefined{KOMAClassName}{% if non-KOMA class
  \IfFileExists{parskip.sty}{%
    \usepackage{parskip}
  }{% else
    \setlength{\parindent}{0pt}
    \setlength{\parskip}{6pt plus 2pt minus 1pt}}
}{% if KOMA class
  \KOMAoptions{parskip=half}}
\makeatother
\usepackage{longtable,booktabs,array}
\newcounter{none} % for unnumbered tables
\usepackage{calc} % for calculating minipage widths
% Correct order of tables after \paragraph or \subparagraph
\usepackage{etoolbox}
\makeatletter
\patchcmd\longtable{\par}{\if@noskipsec\mbox{}\fi\par}{}{}
\makeatother
% Allow footnotes in longtable head/foot
\IfFileExists{footnotehyper.sty}{\usepackage{footnotehyper}}{\usepackage{footnote}}
\makesavenoteenv{longtable}
\setlength{\emergencystretch}{3em} % prevent overfull lines
\providecommand{\tightlist}{%
  \setlength{\itemsep}{0pt}\setlength{\parskip}{0pt}}

\usepackage{booktabs}
\usepackage{array}
\usepackage{textcomp}
\DeclareUnicodeCharacter{2264}{$\leq$}
\DeclareUnicodeCharacter{2265}{$\geq$}
\DeclareUnicodeCharacter{2212}{$-$}
\DeclareUnicodeCharacter{2013}{--}
\DeclareUnicodeCharacter{2014}{---}
\DeclareUnicodeCharacter{03B1}{$\alpha$}
\DeclareUnicodeCharacter{03B2}{$\beta$}
\DeclareUnicodeCharacter{03B3}{$\gamma$}
\DeclareUnicodeCharacter{03B4}{$\delta$}
\DeclareUnicodeCharacter{03B5}{$\epsilon$}
\DeclareUnicodeCharacter{03B7}{$\eta$}
\DeclareUnicodeCharacter{03B8}{$\theta$}
\DeclareUnicodeCharacter{03BB}{$\lambda$}
\DeclareUnicodeCharacter{03BC}{$\mu$}
\DeclareUnicodeCharacter{03BD}{$\nu$}
\DeclareUnicodeCharacter{03C0}{$\pi$}
\DeclareUnicodeCharacter{03C1}{$\rho$}
\DeclareUnicodeCharacter{03C3}{$\sigma$}
\DeclareUnicodeCharacter{03C4}{$\tau$}
\DeclareUnicodeCharacter{03C6}{$\phi$}
\DeclareUnicodeCharacter{03C7}{$\chi$}
\DeclareUnicodeCharacter{03C8}{$\psi$}
\DeclareUnicodeCharacter{03C9}{$\omega$}
\DeclareUnicodeCharacter{0393}{$\Gamma$}
\DeclareUnicodeCharacter{0394}{$\Delta$}
\DeclareUnicodeCharacter{039B}{$\Lambda$}
\DeclareUnicodeCharacter{03A3}{$\Sigma$}
\DeclareUnicodeCharacter{03A9}{$\Omega$}
\DeclareUnicodeCharacter{2202}{$\partial$}
\DeclareUnicodeCharacter{210F}{$\hbar$}
\DeclareUnicodeCharacter{2248}{$\approx$}
\DeclareUnicodeCharacter{2260}{$\neq$}
\DeclareUnicodeCharacter{00B1}{$\pm$}
\DeclareUnicodeCharacter{00D7}{$\times$}
\DeclareUnicodeCharacter{2192}{$\rightarrow$}
\DeclareUnicodeCharacter{221A}{$\sqrt{}$}
\DeclareUnicodeCharacter{221E}{$\infty$}
\DeclareUnicodeCharacter{2272}{$\lesssim$}
\DeclareUnicodeCharacter{2273}{$\gtrsim$}
\DeclareUnicodeCharacter{2261}{$\equiv$}
\DeclareUnicodeCharacter{207A}{$^{+}$}
\DeclareUnicodeCharacter{207B}{$^{-}$}
\DeclareUnicodeCharacter{2070}{$^{0}$}
\DeclareUnicodeCharacter{00B2}{$^{2}$}
\DeclareUnicodeCharacter{00B3}{$^{3}$}
\DeclareUnicodeCharacter{2074}{$^{4}$}
\DeclareUnicodeCharacter{2075}{$^{5}$}
\DeclareUnicodeCharacter{2076}{$^{6}$}
\DeclareUnicodeCharacter{00B9}{$^{1}$}
\DeclareUnicodeCharacter{2077}{$^{7}$}
\DeclareUnicodeCharacter{2078}{$^{8}$}
\DeclareUnicodeCharacter{2079}{$^{9}$}
\DeclareUnicodeCharacter{2080}{$_{0}$}
\DeclareUnicodeCharacter{2081}{$_{1}$}
\DeclareUnicodeCharacter{2082}{$_{2}$}
\DeclareUnicodeCharacter{2083}{$_{3}$}
\DeclareUnicodeCharacter{2084}{$_{4}$}
\DeclareUnicodeCharacter{2089}{$_{9}$}
\renewcommand{\arraystretch}{1.4}
\setlength{\tabcolsep}{8pt}
\usepackage{bookmark}
\IfFileExists{xurl.sty}{\usepackage{xurl}}{} % add URL line breaks if available
\urlstyle{same}
\hypersetup{
  hidelinks,
  pdfcreator={LaTeX via pandoc}}

\author{}
\date{}

\begin{document}

\section{GWT Predictions for the MUSE
Experiment}\label{gwt-predictions-for-the-muse-experiment}

\textbf{Jonathan D. Wollenberg} ORCID:
\href{https://orcid.org/0009-0009-5872-9076}{0009-0009-5872-9076}

March 26, 2026 --- \emph{Prediction documented prior to MUSE
publication}

\begin{center}\rule{0.5\linewidth}{0.5pt}\end{center}

\subsection{Summary of Predictions}\label{summary-of-predictions}

Geometric Wave Theory (GWT) makes the following \textbf{falsifiable
predictions} for the MUSE experiment at PSI, which measures elastic
muon-proton and electron-proton scattering at the same kinematics:

{\def\LTcaptype{none} % do not increment counter
\begin{longtable}[]{@{}
  >{\raggedright\arraybackslash}p{(\linewidth - 4\tabcolsep) * \real{0.3333}}
  >{\raggedright\arraybackslash}p{(\linewidth - 4\tabcolsep) * \real{0.3056}}
  >{\raggedright\arraybackslash}p{(\linewidth - 4\tabcolsep) * \real{0.3611}}@{}}
\toprule\noalign{}
\begin{minipage}[b]{\linewidth}\raggedright
Prediction
\end{minipage} & \begin{minipage}[b]{\linewidth}\raggedright
GWT Value
\end{minipage} & \begin{minipage}[b]{\linewidth}\raggedright
Falsified if
\end{minipage} \\
\midrule\noalign{}
\endhead
\bottomrule\noalign{}
\endlastfoot
Proton radius from μ-p scattering & 0.84062 ± 0.00002 fm & Differs from
e-p by \textgreater{} 0.5\% \\
Proton radius from e-p scattering & 0.84062 ± 0.00002 fm & Differs from
μ-p by \textgreater{} 0.5\% \\
Lepton universality & Exact (at tree level) & Any lepton-specific form
factor \\
Muon charge radius & \textless{} 10⁻²⁰ fm (point-like) & R\_μ
\textgreater{} 0.001 fm \\
Two-photon exchange (μ vs e) & Identical structure &
Lepton-mass-dependent TPE beyond QED \\
\end{longtable}
}

\textbf{Core prediction: MUSE will find no discrepancy between μ-p and
e-p extracted proton radii.} Both will yield r\_p ≈ 0.841 fm, consistent
with muonic hydrogen spectroscopy.

\begin{center}\rule{0.5\linewidth}{0.5pt}\end{center}

\subsection{1. Why GWT Makes These
Predictions}\label{why-gwt-makes-these-predictions}

\subsubsection{1.1 The Proton Radius}\label{the-proton-radius}

In GWT, the proton charge radius is derived from the zero-mode structure
of a sine-Gordon kink on the d=3 cubic lattice:

\[r_p = (d+1)\left(1 - \alpha \cdot \frac{d^3-1}{d^3\pi^2}\right) \frac{\hbar c}{m_p} = 4 \times 0.99981 \times 0.21031 \text{ fm} = 0.84062 \text{ fm}\]

This is a property of the \textbf{proton} (the kink-antikink pair), not
of the scattering probe. The proton looks the same to any lepton because
the proton's structure is topological --- it's determined by the lattice
geometry, not by the probe particle.

\subsubsection{1.2 The Muon Is Point-Like}\label{the-muon-is-point-like}

In GWT, particles are classified by their topological structure on the
lattice:

\begin{itemize}
\tightlist
\item
  \textbf{3D torus (proton, neutron):} Kink-antikink pair wrapping all 3
  spatial dimensions. Has d+1 = 4 zero modes → measurable charge radius
  \textasciitilde{} 0.84 fm.
\item
  \textbf{1D breather (electron):} Oscillation along one lattice
  direction. No topological winding → point-like. Charge radius at
  Planck scale (\textasciitilde10⁻³⁵ m).
\item
  \textbf{Generation excitation (muon):} Same 1D breather structure as
  electron, with mass ratio m\_μ/m\_e = d/((d-1)α) = 205.6. Still 1D →
  still point-like.
\end{itemize}

The muon has \textbf{no topological zero modes} that would generate a
classical charge radius. Its spatial extent is the breather width
\textasciitilde1/ε in Planck units, which is \textasciitilde10⁻³⁵ m ---
seventeen orders of magnitude below MUSE's sensitivity
(\textasciitilde10⁻¹⁸ m).

\subsubsection{1.3 Lepton Universality}\label{lepton-universality}

GWT predicts exact lepton universality at tree level because:

\begin{enumerate}
\def\labelenumi{\arabic{enumi}.}
\tightlist
\item
  The electromagnetic coupling α is derived from the lattice tunneling
  rate, which is the same for all charged particles
\item
  The proton's form factor is a property of the proton (the kink
  topology), not the probe
\item
  The muon and electron differ only in their internal oscillation
  frequency, not in how they couple to the electromagnetic field
\end{enumerate}

Any lepton universality violation in MUSE would require the probe
particle's internal structure to affect the scattering --- which in GWT
would mean the muon has non-trivial topology. This would contradict the
successful derivation of the muon mass as a generation excitation of the
electron.

\begin{center}\rule{0.5\linewidth}{0.5pt}\end{center}

\subsection{2. What MUSE Measures}\label{what-muse-measures}

MUSE (MUon Scattering Experiment) at the Paul Scherrer Institute uses
the πM1 beamline to deliver simultaneous π, μ, e beams at momenta
115--210 MeV/c.~It measures:

\begin{itemize}
\tightlist
\item
  Elastic μ⁺p and μ⁻p differential cross sections
\item
  Elastic e⁺p and e⁻p differential cross sections (same detector, same
  kinematics)
\item
  Ratio of μ-p to e-p cross sections (many systematics cancel)
\end{itemize}

The extracted observable is the proton charge form factor G\_E(Q²) at
low Q² ≈ 0.002--0.08 GeV², from which the charge radius is obtained via:

\[\langle r^2 \rangle = -6 \left.\frac{dG_E}{dQ^2}\right|_{Q^2=0}\]

\subsubsection{2.1 Sensitivity to Muon
Structure}\label{sensitivity-to-muon-structure}

If the muon has a charge radius R\_μ, it modifies the measured cross
section through the muon's own form factor:

\[F_\mu(Q^2) = 1 - \frac{Q^2 R_\mu^2}{6} + \ldots\]

This enters the cross section as a multiplicative factor, effectively
shifting the extracted ``proton radius'' by:

\[\delta r_p^2 \approx R_\mu^2\]

For MUSE's target precision of \textasciitilde1\% on r\_p
(\textasciitilde0.008 fm), this means MUSE is sensitive to R\_μ ≳ 0.008
fm. GWT predicts R\_μ \textasciitilde{} 10⁻²⁰ fm, so \textbf{no muon
structure will be visible}.

\begin{center}\rule{0.5\linewidth}{0.5pt}\end{center}

\subsection{3. Falsification Criteria}\label{falsification-criteria}

GWT would be \textbf{challenged} if MUSE finds:

\begin{enumerate}
\def\labelenumi{\arabic{enumi}.}
\tightlist
\item
  \textbf{r\_p(μ-p) ≠ r\_p(e-p)} at \textgreater{} 3σ significance

  \begin{itemize}
  \tightlist
  \item
    This would imply lepton non-universality or muon structure
  \item
    GWT has no mechanism for this (both probe the same kink topology)
  \end{itemize}
\item
  \textbf{r\_p ≠ 0.841 fm} from either channel

  \begin{itemize}
  \tightlist
  \item
    The GWT value is 0.84062 fm with \textasciitilde0.02\% theoretical
    uncertainty
  \item
    A measurement of 0.87 fm (the old CODATA value) would be problematic
  \end{itemize}
\item
  \textbf{Anomalous Q²-dependence} in μ-p vs e-p ratio

  \begin{itemize}
  \tightlist
  \item
    A flat ratio = lepton universality (GWT prediction)
  \item
    A Q²-dependent ratio = new physics beyond GWT
  \end{itemize}
\end{enumerate}

GWT would be \textbf{confirmed} if MUSE finds:

\begin{enumerate}
\def\labelenumi{\arabic{enumi}.}
\tightlist
\item
  r\_p(μ-p) = r\_p(e-p) = 0.841 ± 0.01 fm (consistent with muonic H)
\item
  No lepton universality violation at any Q²
\item
  Muon consistent with point-like (no anomalous form factor)
\end{enumerate}

\begin{center}\rule{0.5\linewidth}{0.5pt}\end{center}

\subsection{4. Connection to the Proton Radius
Puzzle}\label{connection-to-the-proton-radius-puzzle}

The proton radius puzzle (2010--2019) arose from a discrepancy between:
- Muonic hydrogen spectroscopy: r\_p = 0.84087 ± 0.00039 fm - Electronic
hydrogen + e-p scattering: r\_p = 0.8751 ± 0.0061 fm

GWT predicted the muonic value from the start (the formula gives 0.84062
fm). The resolution of the puzzle in favor of \textasciitilde0.841 fm is
consistent with GWT.

MUSE provides the \textbf{definitive test}: same detector, same
kinematics, both leptons. If μ-p and e-p give the same r\_p ≈ 0.841 fm,
the puzzle is closed and GWT's prediction is confirmed.

\begin{center}\rule{0.5\linewidth}{0.5pt}\end{center}

\subsection{5. The Deeper Question: Why Is the Muon
Point-Like?}\label{the-deeper-question-why-is-the-muon-point-like}

In the Standard Model, the muon is point-like by assumption (it's a
fundamental fermion). GWT provides a \textbf{reason}: the muon is a 1D
oscillation (breather) on the lattice, with no topological winding. Only
3D topological objects (kink-antikink tori) have measurable spatial
extent.

This connects to the correction hierarchy discovered in GWT: -
\textbf{1D/2D modes (leptons):} No lattice distortion → no corrections
to bare values → point-like - \textbf{3D torus (proton):} Wraps all 3
dimensions → local lattice distortion → effective c reduced → correction
terms (1 - α·26/(27π²)) → measurable radius

The muon mass formula m\_μ/m\_e = d/((d-1)α) = 205.6 (obs: 206.8, 0.6\%)
works \textbf{without} any spatial correction, confirming the muon has
no 3D topological structure.

Supporting evidence: all three lepton masses (e, μ, τ) are predicted
without corrections: - m\_μ/m\_e = d/((d-1)α) = 205.6 (0.6\%) - Koide
parameter = (d-1)/d = 2/3 (8.8 ppm) - m\_τ from Koide = 1777 MeV
(0.006\%)

\begin{center}\rule{0.5\linewidth}{0.5pt}\end{center}

\subsection{6. Muon g-2: A Derived
Prediction}\label{muon-g-2-a-derived-prediction}

\textbf{Update (March 28, 2026): The muon g-2 has been fully derived.}

The electron g-2 uses the full \(O_h\) channel decomposition:

\[a_e = \frac{\alpha}{2\pi}\left(1 - \frac{\alpha}{2d-1} - \frac{\alpha^2}{2d+1}\right) = 0.00115965182 \quad (\text{obs: ...218, } {-0.31 \text{ ppm}})\]

The muon g-2 adds a hadronic VP correction with an NLO term from the
\(O_h \to D_{4h}\) subgroup restriction:

\[a_\mu = a_e + \frac{\alpha^2}{2\pi}\left(\frac{m_\mu}{m_\pi}\right)^2 \frac{d}{d-1} \cdot F_{O_h} \cdot F_{D_{4h}}\]

where: -
\(F_{O_h} = \frac{169}{198} = \frac{(d^2+d+1)^2}{2d^2(d^2+d-1)}\) ---
bifundamental EM×QCD trace (Oh symmetry of 3D vacuum) -
\(F_{D_{4h}} = 1 + \alpha \cdot \frac{d^2+d-1}{d^2+1} = 1 + \alpha \cdot \frac{11}{10}\)
--- NLO from D4h restriction

\textbf{Result:} \(a_\mu = 0.00116591962\) (obs: \(0.00116592061\),
\textbf{−0.85 ppm})

The SM predicts \(0.00116591810\) (−2.15 ppm). \textbf{GWT is 2.5×
closer to observation.}

\textbf{Derivation:} The muon lives on \(d-1 = 2\) axes of the \(d=3\)
cube. Its local symmetry is \(D_{4h}\) (the square), not \(O_h\) (the
cube). Restricting \(T_{1u} \otimes T_{1u}\) from \(O_h\) to \(D_{4h}\):

\[T_{1u} \to A_{2u} + E_u \quad \text{in } D_{4h}\]

\[(A_{2u} + E_u) \otimes (A_{2u} + E_u) \text{ has } \textbf{2} \text{ } A_{1g} \text{ channels (not 1)}\]

The extra \(A_{1g}\) comes from the \(E_g\) irrep of \(O_h\) splitting
under \(D_{4h}\): \(E_g \to A_{1g} + B_{1g}\). The \(d_{z^2}\) component
looks scalar from the 2D muon frame. This channel is
\(\alpha\)-suppressed and normalized by the ratio of QCD exchange paths
to coupling modes: \((d^2+d-1)/(d^2+1) = 11/10\).

\textbf{Key identity:} \(d^2+d-1 = d^2+2\) only at \(d=3\) (requires
\(d-1=2\)).

The D4h character table, branching rules, and full calculation are in
\texttt{reference/nuclear.md} and
\texttt{calculations/core/muon\_g2\_d4h.py}.

\begin{center}\rule{0.5\linewidth}{0.5pt}\end{center}

\subsection{7. Timeline and Experimental
Status}\label{timeline-and-experimental-status}

\begin{itemize}
\tightlist
\item
  \textbf{MUSE commissioning:} 2019--2022
\item
  \textbf{Production data:} 2023--2025
\item
  \textbf{First results:} Expected 2025--2026
\item
  \textbf{This prediction documented:} March 26, 2026
\end{itemize}

All GWT predictions in this document were derived from the Lagrangian
\(\mathcal{L} = \frac{1}{2}(\partial\phi)^2 + \frac{1}{\pi^2}(1-\cos\pi\phi)\)
on the d=3 cubic lattice with zero free parameters.

\begin{center}\rule{0.5\linewidth}{0.5pt}\end{center}

\subsection{References}\label{references}

{[}1{]} Antognini et al., Science 339, 417 (2013) --- Muonic hydrogen
r\_p = 0.84087 fm {[}2{]} Mohr et al., Rev.~Mod. Phys. 84, 1527 (2012)
--- CODATA 2010 r\_p = 0.8775 fm {[}3{]} Brandt et al., Phys. Rev.~Lett.
(2026) --- PRad-II r\_p = 0.840615 fm (to be confirmed) {[}4{]}
Alexandrou et al., Phys. Rev.~D 101, 114513 (2020) --- Lattice QCD r\_p
{[}5{]} Wollenberg, ``Proton-Electron Mass Ratio from Sine-Gordon
Lattice'' (2026)

\end{document}
